As conclusion, we can say that the numbers of bleedings and reheatings have a significant impact on the cycle's efficiency.
These numbers can vary depending on the application and the constraints. We also see that the Rankine cycle can be improved 
by using supercritical pressures. Compared to a gas power plant, the efficiency is higher but the constraints are more important.
Indeed, we need a source of cooling fluid and so, another cycle designed to manage the heat discharged by the condenser. Even though the 
Rankine cycle has been used for a long time with non-renewable sources, it can still be used in more modern applications. We can think
of combined cycle power plants or ROC's (Rankine Organic Cycles) which create electricity from heat waste or renewable sources.